\section{LlamaIndex：构建多模态知识助手}

\subsection{从 Basic RAG 到 Advanced RAG}

Jerry Liu 指出 Basic RAG 的四大局限：
\begin{enumerate}
    \item \textbf{数据处理粗糙}：简单的文本分块无法处理表格、图片、复杂布局
    \item \textbf{LLM 仅用于合成}：没有利用推理和规划能力
    \item \textbf{无状态}：每次查询独立，无个性化记忆
    \item \textbf{输出单一}：只能生成简单文本回答
\end{enumerate}

\begin{warningbox}{Garbage In Garbage Out}
RAG 应用的质量取决于数据处理管线的质量。如果 PDF 解析器无法正确提取表格和图片，后续的 LLM 推理无论多强大都无法弥补。数据质量是生产级 LLM 应用的必要条件。
\end{warningbox}

\subsection{多模态 RAG 管线}

构建多模态 RAG 的关键步骤：

\begin{enumerate}
    \item \textbf{高质量文档解析}：将 PDF 中的文本、表格、图表、图片分别提取为结构化元素
    \item \textbf{层次化索引}：为每个元素生成多种文本表示（摘要、子块等），指向源元素
    \item \textbf{多模态检索}：给定查询，检索文本和图片 chunk，利用多模态 LLM 同时处理文字和视觉信息
\end{enumerate}

\begin{knowledgebox}{多模态模型的新机遇}
当前主流模型（GPT-4o、Claude 3.5 Sonnet、Gemini）都支持同时输入文字和图片。这意味着即使使用文本嵌入模型建立索引，在最终推理阶段仍可以将原始图片喂入 LLM，实现真正的视觉理解。
\end{knowledgebox}

\subsection{Agentic RAG：在 RAG 之上添加推理层}

\begin{importantbox}{Agentic RAG 的核心思想}
将 RAG 检索视为一种工具（tool），在其上层添加 Agent 推理层。Agent 可以决定如何分解问题、调用哪些检索工具、何时进行反思验证，从而处理基础 RAG 无法回答的复杂问题（摘要、多文档比较、多步骤分析等）。
\end{importantbox}

两种架构选择：
\begin{itemize}
    \item \textbf{约束型流程}：开发者手动定义 if-else 逻辑和路由规则，更可靠但灵活性低
    \item \textbf{通用 Agent 架构}（如 ReAct）：Agent 自主选择工具和推理路径，更灵活但可靠性较低、成本较高
\end{itemize}

经验法则：通用 Agent 架构中，工具数量尽量控制在 4--5 个，不超过 10 个。

\subsection{生产部署}

LlamaIndex 提供事件驱动的编排系统，支持：
\begin{itemize}
    \item 将每个 Agent 工作流封装为微服务 API
    \item 通过中央消息队列实现 Agent 间通信
    \item Human-in-the-loop：Agent 可暂停等待人类输入后继续执行
    \item 支持本地开发和 Kubernetes 部署
\end{itemize}

\subsection{本章小结}

LlamaIndex 从数据处理质量出发，构建了从文档解析到多模态检索再到 Agentic 推理的完整管线。其核心理念是：好的数据质量是基础，Agent 推理是上层建筑，两者缺一不可。

